\documentclass[11pt]{article}
\usepackage{wsi-style-lua}
\usepackage{amsmath}
\newcommand{\SC}{S_{\mathrm{C8}}}
\newcommand{\Sesc}{S_{\mathrm{esc}}}
\wsirunhead{On the Status of Local Entropy-Current Extensions}
\title{On the Status of Local Entropy-Current Extensions of the\\[2pt]
\large Gravitational Entropy Escrow Framework: A Clarification Note}
\author{Grant Whitmer}
\date{May 2026}

\begin{document}
\maketitle

\begin{abstract}
In dialogue exploring covariant extensions of the Gravitational Entropy Escrow framework
(Whitmer 2026), several large-language-model systems independently proposed an equation of the form
$\nabla_{\mu} J_{S}^{\mu} = (2\pi k_{B}/\hbar)\,\varepsilon_{\mathrm{bind}}$, designated~C8, in which a
local entropy 4-current is sourced by the local energy density. We document its resolution. With
dimensions accounted for correctly, the standard cosmological-constant energy density
$\varepsilon_{\Lambda} = \Lambda c^{4}/(8\pi G)$, and the integration time chosen as the horizon
light-crossing time, C8 reproduces both the Bekenstein--Hawking entropy of a Schwarzschild black hole
and the Gibbons--Hawking entropy of the de Sitter horizon exactly. This recovery is a consistency
check, not a derivation: C8 with $t = R/c$ is algebraically identical to the saturated Bekenstein
bound (Bekenstein 1981), so any horizon that saturates that bound is reproduced trivially. The choice
$t = R/c$ is post-hoc---other natural horizon time scales give answers spanning many orders of
magnitude. C8 therefore reduces to a rate-equation reformulation of a 1981 result rather than new
physics. The published framework's interpretive content is unaffected, and its central open problem
(the order-unity coefficient $\alpha \approx 1.34$ between the bare $\Lambda$-form $a_{0}$ and the
empirical Milgrom value) is untouched by C8.
\end{abstract}

\section{Motivation}

The Gravitational Entropy Escrow framework (Whitmer 2026, henceforth W26) proposes that gravitational
binding energy is structurally identical to entropy held in escrow against the local Unruh
temperature: $\Sesc = |U_{\mathrm{grav}}|/T$, with the force law $F = T\,\nabla\Sesc$ recovering
Newton's law (Verlinde 2011). The framework recovers Bekenstein--Hawking entropy at Schwarzschild
horizons, the order of magnitude of the empirical Milgrom acceleration scale $a_{0}$ from a
thermodynamic crossover at the de Sitter horizon, and several other results discussed in W26.

In dialogue with multiple large-language-model systems exploring such an extension, a covariant
equation of the form
\begin{equation}
  \nabla_{\mu} J_{S}^{\mu} = \frac{2\pi k_{B}}{\hbar}\,\varepsilon_{\mathrm{bind}}
  \tag{C8}
\end{equation}
was independently proposed, with $J_{S}^{\mu}$ an entropy 4-current and $\varepsilon_{\mathrm{bind}}$
the local energy density. The dimensions balance: both sides have units of entropy-production density
per unit time, $k_{B}/(\mathrm{m}^{3}\,\mathrm{s})$. This note documents the resolution of C8's status.
The result is small and specific: C8 with the light-crossing-time integration prescription is
algebraically identical to the saturated Bekenstein bound (Bekenstein 1981), so it reproduces both the
Bekenstein--Hawking and Gibbons--Hawking entropies exactly when the standard cosmological-constant
energy density is used. The recovery is a consistency check rather than a derivation, because horizons
saturate the Bekenstein bound by construction, and the choice of integration time is post-hoc rather
than derived.

\section{C8 as a rate equation}

C8 is a continuity-like equation in which the divergence of an entropy 4-current is sourced by an
energy density; both sides have dimensions of entropy-production density per unit time. To compare C8
to a total entropy (such as Bekenstein--Hawking or Gibbons--Hawking) one must first integrate the rate
over both volume and time. Integrating over a region of volume $V$ containing total binding energy $E$
gives the total entropy-production rate, and multiplying by a characteristic time $t_{h}$ yields a
total entropy:
\begin{equation}
  \frac{dS}{dt} = \frac{2\pi k_{B}}{\hbar}\,E,
  \qquad
  \SC = \frac{2\pi k_{B}}{\hbar}\,E\,t_{h}.
\end{equation}
The choice of $t_{h}$ is the load-bearing question; it is not specified by C8 itself.

\section{C8 with $t = R/c$ is the saturated Bekenstein bound}

Choosing $t_{h} = R/c$, where $R$ is the radius of the region:
\begin{equation}
  \SC = \frac{2\pi k_{B}}{\hbar}\,E\,\frac{R}{c} = \frac{2\pi k_{B} E R}{\hbar c}.
\end{equation}
This is the Bekenstein bound (Bekenstein 1981), saturated. The bound applies as an inequality to any
localized system: the maximum entropy a region of energy $E$ and radius $R$ can contain in a quantum
field theory respecting causality and the second law is $2\pi k_{B} E R/(\hbar c)$. Black-hole event
horizons saturate the bound, with the saturation value being the Bekenstein--Hawking entropy; the de
Sitter cosmological horizon also saturates it when the appropriate energy is the total vacuum energy
enclosed. Therefore C8 with light-crossing-time integration equals the saturated Bekenstein bound by
algebra---a known relation in a different form, not new physics.

\section{Recovery of horizon entropies}
\subsection{Schwarzschild black hole}
For a Schwarzschild black hole of mass $M$, taking the rest energy $U = Mc^{2}$ and the Schwarzschild
radius $r_{s} = 2GM/c^{2}$ as the relevant radius:
\begin{equation}
  \SC(\mathrm{BH}) = \frac{2\pi k_{B}}{\hbar}\,(Mc^{2})\,\frac{r_{s}}{c}
  = \frac{4\pi G M^{2} k_{B}}{\hbar c} = S_{\mathrm{BH}}.
\end{equation}
The ratio $\SC(\mathrm{BH})/S_{\mathrm{BH}} = 1$ exactly, independent of $M$, verified numerically and
symbolically in the accompanying script.

\subsection{de Sitter cosmological horizon}
For the de Sitter horizon, the radius is $R_{\mathrm{dS}} = \sqrt{3/\Lambda}$. Using the standard
energy density of the cosmological constant,
\begin{equation}
  \varepsilon_{\Lambda} = \frac{\Lambda c^{4}}{8\pi G} \approx 5.3\times 10^{-10}\ \mathrm{J/m^{3}}
  \quad(\text{Planck 2018}),
\end{equation}
the closely related quantity $\Lambda c^{2}/(8\pi G)$ is the equivalent \emph{mass} density
($\approx 6.0\times 10^{-27}\,\mathrm{kg/m^{3}}$) and must not be substituted for the energy density.
With total enclosed vacuum energy $E_{\mathrm{total}} = \varepsilon_{\Lambda} V_{\mathrm{dS}}$,
$V_{\mathrm{dS}} = (4\pi/3)R_{\mathrm{dS}}^{3}$,
\begin{equation}
  \SC(\mathrm{dS}) = \frac{2\pi k_{B}}{\hbar}\,(\varepsilon_{\Lambda} V_{\mathrm{dS}})\,
  \frac{R_{\mathrm{dS}}}{c} = \frac{3\pi c^{3} k_{B}}{G \hbar \Lambda} = S_{\mathrm{GH}}.
\end{equation}
The ratio $\SC(\mathrm{dS})/S_{\mathrm{GH}} = 1$ exactly, independent of $\Lambda$.

\section{The light-crossing time is post-hoc}

The choice $t_{h} = R/c$ is not derived from C8; it is one of several time scales naturally associated
with a horizon. For a solar-mass Schwarzschild black hole, alternative natural choices include:
\begin{itemize}
\item Light-crossing time $r_{s}/c \approx 10^{-5}$~s, giving ratio $\SC/S_{\mathrm{BH}} = 1$.
\item Inverse surface gravity $c/\kappa \approx 2\times 10^{-5}$~s, giving ratio $= 2$.
\item $\pi r_{s}/c$ (inverse Hawking frequency) $\approx 3\times 10^{-5}$~s, giving ratio $= \pi$.
\item Hawking evaporation time $\tau_{\mathrm{evap}} \approx 2\times 10^{67}$~yr, giving ratio
$\approx 10^{80}$.
\end{itemize}
Only the light-crossing time produces ratio $= 1$, and it does so because $R/c$ is precisely the
factor needed for C8 with prefactor $2\pi k_{B}/\hbar$ to algebraically equal the Bekenstein bound.
The other natural time scales---each defensible on its own physical grounds---give different answers
across many orders of magnitude. This is the textbook signature of a post-hoc choice: a genuine
derivation would either select $t_{h}$ from a deeper principle or survive without privileging any
particular choice. The agreement with horizon entropies that $R/c$ produces is therefore a consequence
of the algebraic identity, not independent evidence for C8.

\section{Status and implications}

C8 with the light-crossing-time prescription is algebraically equivalent to the saturated Bekenstein
bound; it contains no information beyond Bekenstein 1981 plus the choice of integration time. The
published escrow framework (W26) is unaffected: the static postulate $\Sesc = |U_{\mathrm{grav}}|/T$,
the recovery of Newton's law, the recovery of Bekenstein--Hawking entropy at horizons (W26 \S5), the
identification of the deep-MOND limit with a de Sitter thermodynamic crossover (W26 \S6), and the
empirical tests of W26 \S\S6.1, 8.4--8.6---the Genzel five-case test, the SPARC reanalysis, the baryonic Tully--Fisher relation, and cross-redshift consistency (all of which depend on the deep-MOND asymptote $a_{0} = c^{2}\sqrt{\Lambda/3}/(2\pi)$, not on C8)---all stand.

What is ruled out is the natural-looking local-current generalization of the static escrow postulate
as a candidate independent extension. Three directions remain open and are \emph{not} ruled out. First, \emph{quantum corrections to horizon temperatures}: the leading-order Gibbons--Hawking and Hawking temperatures may receive corrections from entanglement-entropy structure across horizons, depending on the cosmological matter content ($\Omega_m$); such corrections might derive the order-unity coefficient $\alpha \approx 1.34$ named as the open problem in W26 \S8.4.3. Calculations of this type exist for lower-dimensional analogs (Calabrese--Cardy 2004) but not, to our knowledge, for the four-dimensional de Sitter horizon in a universe with realistic matter content. Second, \emph{Rindler--de Sitter horizon matching}: the framework's commitment to a constant $a_0$ requires the de Sitter horizon temperature to set the floor for entropy exchange of accelerated observers; a correct treatment of the matching between local Rindler horizons and the global de Sitter horizon would not be a local continuity equation. Third, \emph{lattice quantum-field-theory tests of the static escrow postulate}: the identification $\Sesc = |U_{\mathrm{grav}}|/T$ can be tested directly by computing the entanglement entropy across a boundary separating two localized mass distributions and checking whether it scales as $|U_{\mathrm{grav}}|/T_{\mathrm{Unruh}}$---computationally tractable in $1{+}1$D using Peschel's formula for Gaussian states (Peschel 2003), and a sharp falsification test of the static postulate. The
framework's central open derivation problem ($\alpha \approx 1.34$) lives in the deep-MOND
interpolation, not in horizon-entropy magnitudes, and C8 does not address it in either direction.

\section{Methodological notes}

This paper is the third draft of an analysis of C8; two earlier drafts reported wrong results for
opposite reasons, documented here because the failure modes are easy to repeat.

The first draft (v0.1) reported that C8 fails by 30 orders of magnitude for galactic systems---a
dimensional category error in which the C8 source term (an entropy-production-rate density) was
compared directly to total entropies. It would have been caught immediately by a symbolic dimensional
check, which was not performed on the first pass. The second draft (v0.2) corrected that error but
reported that C8 underpredicts Gibbons--Hawking by $1/c^{2}$; this used the mass density
$\Lambda c^{2}/(8\pi G)$ where the energy density $\Lambda c^{4}/(8\pi G)$ was required. Symbolic
dimensional analysis cannot catch this, because both forms are dimensionally consistent within their
respective conventions; only physical identification (``this should be an energy density, not a mass
density'') catches it. The third draft (this document) compares each computed quantity against a
published reality-check value---the Planck 2018 dark-energy density, the de Sitter horizon radius, and
the Gibbons--Hawking entropy---before treating it as established.

Notably, after one LLM audit identified the mass-vs-energy density error, an ``audit of the audit''
from another system confidently reversed the correct result. Resolving the disagreement required
independent first-principles calculation against three external anchors: the SI unit definitions, the
Planck 2018 values, and the Friedmann equation $H^{2} = \Lambda c^{2}/3$. Three methodological lessons
follow: symbolic dimensional analysis is necessary but not sufficient (it checks consistency within a
convention, not the choice of convention); reality checks against published values are necessary; and
multi-LLM adversarial review is load-bearing but cannot resolve disagreements among LLMs---recursion of
LLM review without external grounding is unstable.

\section{Conclusion}

C8 is the saturated Bekenstein bound (Bekenstein 1981) expressed as a local rate equation. With the
light-crossing-time integration prescription and the standard cosmological-constant energy density, it
reproduces Bekenstein--Hawking and Gibbons--Hawking entropies exactly---a consistency check rather than
a derivation, because both horizons saturate the Bekenstein bound by construction, and the choice of
integration time is post-hoc. C8 therefore does not constitute an independent covariant extension of
the escrow framework, even though it is internally consistent. The published framework (Whitmer 2026)
is unaffected; its central open problem ($\alpha \approx 1.34$) is not addressed. Future researchers
should pursue extensions through quantum corrections to horizon temperatures, Rindler--de Sitter
matching, or lattice QFT tests of the static postulate, none of which are ruled out here.

\subsection*{Acknowledgments}
The candidate equation C8 was first proposed in dialogue with several large-language-model systems, which independently proposed equivalent forms. Adversarial review across multiple rounds was contributed by these systems; a mass-versus-energy density convention error in an intermediate draft was identified by one system's audit, and a subsequent ``audit of the audit'' from another confidently reversed the correction. Resolving the disagreement required independent first-principles calculation against Planck 2018 published values. The author is solely responsible for the calculations, conclusions, and any remaining errors.

\subsection*{Code and data availability}
The Python script reproducing all numerical claims (\texttt{c8\_calculations\_v0\_3.py}) is deposited
alongside this preprint. It uses numerical evaluation and symbolic dimensional analysis (sympy) and
includes explicit reality checks against published values for the dark-energy density, the de Sitter
horizon radius, and the Gibbons--Hawking entropy. It runs in under one second and is licensed MIT.

\subsection*{Data, Code \& Resources}
\noindent Manuscript source and archived record: \url{https://github.com/Windstorm-Institute/c8-clarification-note} (Zenodo: \href{https://doi.org/10.5281/zenodo.20041991}{10.5281/zenodo.20041991}). Experiments, datasets, and analysis code: \url{https://github.com/Windstorm-Labs/c8-clarification-note}. Windstorm Institute: \url{https://windstorminstitute.org}. Windstorm Labs: \url{https://windstormlabs.com}.

\begin{thebibliography}{99}\small
\bibitem{bek} J.~D. Bekenstein, \textit{Universal upper bound on the entropy-to-energy ratio for
bounded systems}, Phys.\ Rev.\ D \textbf{23}, 287 (1981).
\bibitem{cc} P.~Calabrese and J.~Cardy, \textit{Entanglement entropy and quantum field theory}, J.\
Stat.\ Mech.\ P06002 (2004).
\bibitem{gh} G.~W. Gibbons and S.~W. Hawking, \textit{Cosmological event horizons, thermodynamics, and
particle creation}, Phys.\ Rev.\ D \textbf{15}, 2738 (1977).
\bibitem{hawking} S.~W. Hawking, \textit{Particle creation by black holes}, Commun.\ Math.\ Phys.\
\textbf{43}, 199 (1975).
\bibitem{padma} T.~Padmanabhan, \textit{Thermodynamical aspects of gravity: new insights}, Rep.\
Prog.\ Phys.\ \textbf{73}, 046901 (2010). arXiv:0911.5004.
\bibitem{peschel} I.~Peschel, \textit{Calculation of reduced density matrices from correlation
functions}, J.\ Phys.\ A \textbf{36}, L205 (2003).
\bibitem{planck} Planck Collaboration, \textit{Planck 2018 results.\ VI.\ Cosmological parameters},
Astron.\ Astrophys.\ \textbf{641}, A6 (2020).
\bibitem{verlinde} E.~P. Verlinde, \textit{On the origin of gravity and the laws of Newton}, JHEP
\textbf{04}, 029 (2011).
\bibitem{w26} G.~L. Whitmer, \textit{Gravitational Entropy Escrow: An Interpretive Synthesis of
Thermodynamic Approaches to Gravity}, Zenodo (2026). DOI:10.5281/zenodo.20031931.
\end{thebibliography}

\end{document}
